\section{Introduction}
\label{sec:introduction}

Low-cost air sensors record particulate matter almost continuously, whereas
regulatory or research-grade instruments operate at selected sites and times.
For exposure assessment, the distribution of the reference measurement at a
given humidity can be more informative than its conditional mean: upper-tail
probabilities and quantiles describe episodes that an average conceals. Sensor
readings are plentiful, but a calibration relation that is adequate overall
may be inaccurate at that humidity or in that tail.

Pulse oximetry, expert review, and delayed-outcome studies create the same
statistical problem. Covariates \(X\) and a surrogate \(S\) are observed for
most units, while the reference outcome \(Y\) is observed only in a verified
subset. We seek the conditional distribution of \(Y\) at \(X=x_0\). The
effective reference sample is therefore the verified sample near \(x_0\),
which may be small even in a large study.

Verified outcomes support a natural estimator that uses an \(X\)-based outcome
regression but leaves the surrogate out of the augmentation. We use this
estimator as the anchor. Adding \(S\) to the regression can lower variance when
it explains the threshold indicators for \(Y\). The same addition can increase
variance when the fitted surrogate correction is poor in the target
neighborhood. The relevant comparison is the local variance change induced by
the correction. A useful method must estimate how much of the correction to
retain for the distribution at \(x_0\).

The CDF target makes this comparison more demanding than a scalar estimation
problem. Threshold-by-threshold tuning can produce unrelated borrowing
decisions along the same distribution. Tuning on the observations used to
evaluate the final CDF also complicates a simultaneous confidence band. We
therefore need one borrowing decision across the threshold range and an
evaluation sample that has played no part in making that decision.

Validation-sample and two-phase methods recover information about variables
measured in a smaller second-phase sample
\citep{Pepe1992,TaoZengLin2017}. Augmented inverse-probability methods provide
the corresponding missing-data representation and efficiency benchmark
\citep{RobinsRotnitzkyZhao1994,Tsiatis2006,CaoTsiatisDavidian2009}. The data
layout can also be viewed as semi-supervised: verified \(Y\) values act as
labels, and the remaining units provide \(X\) and \(S\). Much of the earlier
semi-supervised and prediction-powered literature targets prediction or a
finite-dimensional population parameter
\citep{ChakraborttyCai2018,ZhangBrownCai2019,AzrielEtAl2022,AngelopoulosEtAl2023,AngelopoulosDuchiZrnic2023,ZrnicCandes2024,LiIgnatiadis2025,MiaoEtAl2025}.
General semi-supervised theory already shows how auxiliary information can be
weighted without imposing full borrowing. \citet{SongLinZhou2024} derive
projection-based optimal weights for general M-estimation and cross-validated
inference, while \citet{DengEtAl2024} construct optimal and safe estimators for
high-dimensional regression under possible model misspecification.
\citet{ChenEtAl2025} combine surrogate augmentation with regularized weighting
and adaptive validation sampling for finite-dimensional estimating equations.
Related work examines surrogate validity and efficiency in other inferential
settings \citep{JiLeiZrnic2025,KallusMao2025}.
These procedures choose or regularize auxiliary-information use for parameter
estimation rather than for a threshold-indexed local distribution.

\citet{SuiEtAl2026} localize prediction-powered inference for scalar conditional functionals at
\(x_0\). \citet{TakeishiDingHe2026} use convolution smoothing to develop
prediction-powered inference for finite-dimensional linear quantile-regression
coefficients under possible misspecification, using separate labeled and
unlabeled samples with a widely observed surrogate. \citet{TestaEtAl2025}
allow missing-at-random labels and decreasing
overlap for general smooth semiparametric targets. \citet{Wang2026} develops
uniform inference for transported counterfactual CDFs and quantiles with
surrogate-assisted missing primary outcomes, whereas local
regression-distribution methods assume that the outcome is fully observed
\citep{CattaneoJanssonMa2024,CattaneoEtAl2024}. Here \(X\) and \(S\) are
broadly observed, \(Y\) is selectively verified with \(p_M\) possibly tending
to zero, and the target is the full process \(F_0(\cdot\mid x_0)\). The
effective reference information is \(Mp_Mh^d\); the adaptive decision is one
weight on an anchor-to-surrogate correction shared across thresholds; and the
output comprises a uniform CDF process, simultaneous bands, and conditional
quantile inversion. This combination creates a local, target-specific
negative-transfer problem.

Finite-dimensional optimal or safe weights do not supply a shared path-oracle
result when localization and verification both vanish, nor do scalar
conditional procedures provide one borrowing decision across a CDF followed by
a post-selection simultaneous band.
Theorems~\ref{thm:oracle-equivalence}--\ref{thm:bootstrap} establish these
results for the present data structure.

Our construction uses a property of augmented inverse-probability estimation:
with the correct verification probability, the outcome regression affects
variance without changing the target. We fit an anchor regression using \(X\)
and a second regression using \((X,S)\). Their difference defines a candidate
surrogate correction. One weight, shared across thresholds, is selected by an
observable local variance criterion. A weight of zero gives the anchor and a
weight of one gives full borrowing. Separate samples fit the regressions,
select the weight, and evaluate the CDF; rotating these roles lets every
observation contribute to evaluation.

We introduce an adaptive borrowing rule for a local
reference-outcome CDF: a single estimated weight governs the surrogate
correction across thresholds while evaluation remains independent of fitting
and selection. The theory establishes path-oracle equivalence, a uniform
Gaussian limit and multiplier bands when the verification fraction may
decrease, together with bias conditions for point-target CDF and quantile
inference. Simulations test the transfer criterion under aligned and locally
reversed surrogate relations, and the EPA, VitalDB, and AI-assisted review
studies examine local distribution recovery under controlled reference
sampling. Section~\ref{sec:method} develops the estimator,
Section~\ref{sec:theory} gives the large-sample results,
Section~\ref{sec:simulation} reports the simulations,
Section~\ref{sec:applications} presents two empirical applications and a
controlled review benchmark, and
Section~\ref{sec:discussion} concludes.
